\begin{figure}[!ht]
\centering
\begin{tikzpicture}[x=1cm,y=1cm,>=Stealth,
  vertex/.style={circle,draw,fill=white,inner sep=2pt},
  every node/.style={font=\small}]
\node[vertex] (v) at (0,0) {$v$};
\node[vertex] (b) at (-1.8,-1.0) {$b$};
\node[vertex] (c) at (1.8,-1.0) {$c$};
\node[vertex] (b1) at (-2.6,-2.0) {$b_1$};
\node[vertex] (b2) at (-1.0,-2.0) {$b_2$};
\draw[->,thick,blue!65!black] (v)--node[above left] {$A_v(b)$} (b);
\draw[->] (v)--node[above right] {$A_v(c)$} (c);
\draw[->] (b)--(b1);
\draw[->] (b)--(b2);
\draw[densely dotted] (b1)--(-2.9,-2.6);
\draw[densely dotted] (b1)--(-2.3,-2.6);
\draw[densely dotted] (b2)--(-1.3,-2.6);
\draw[densely dotted] (b2)--(-0.7,-2.6);
\node[align=left,anchor=west] at (0.2,-2.0)
  {Every label below $b$ avoids $A_v(b)$.\\
   Hence $D_b(W)$ uses only $W\setminus A_v(b)$.};
\node[anchor=south] at (0,0.4) {The same set $W$ is used at every node.};
\end{tikzpicture}
\caption{Disjointness along paths controls the change of measure. Every
arc below $b$ lies on a path through its incoming arc, so the potential
$D_b$ is a function of coordinates outside the incoming label. The diagram
is schematic; labels on different branches may overlap.}
\label{fig:outside}
\end{figure}
